\exercisegroup

\begin{enumerate}
\def\labelenumi{\arabic{enumi})}
\item
  设 I 为不可数集，E 为赋有 I, p.~24, 练习 14 中所定义的拓扑的空间 \(\mathbf{R}^{(I)}\)，设 \(E'\) 为其对偶（IV, p.~50, 练习 16, c)）。证明在 \(E'\) 中存在（对 \(\sigma(E', E)\) 而言）非相对紧的子集 H，使得从 H 中点的每个序列都能抽出一个对 \(\sigma(E', E)\) 收敛于 H 中一点的序列。
\item
  \begin{enumerate}
  \def\labelenumii{\alph{enumii})}
  \tightlist
  \item
    设 X 为正则空间，A 为 X 的子集。假设 A 中点的每个序列在 X 中都有一个极限点，并且 X 上存在一个比给定拓扑 \(T_{0}\) 更粗的可度量化拓扑 T。证明 A 在 X 中的闭包 \(\overline{A}\) 是紧可度量化空间（用反证法证明 T 与 \(T_{0}\) 在 \(\overline{A}\) 上诱导的拓扑相同）。
  \end{enumerate}
\end{enumerate}

\begin{enumerate}
\def\labelenumi{\alph{enumi})}
\setcounter{enumi}{1}
\item
  设 E 为分离的局部凸空间，它是一列 \((\mathrm{E}_{n})\) 的并，诸向量子空间对由 E 的拓扑诱导的拓扑可度量化。证明：为使 E 的子集 A 在 E 中相对紧，必要且充分的是 A 中点的每个序列在 E 中有一个极限点。（若 \((\mathrm{W}_{mn})\)（对 \(m \geqslant 1\)）是 \(E_{n}\) 中由在 \(E_{n}\) 中为开的凸零邻域组成的基本邻域系，令 \(U_{mn}\) 为 E 中满足 \(E_{n} \cap U_{mn} = W_{mn}\) 的凸零邻域（II, p.~33, 引理 2）；在 E 上考虑使诸 \(U_{mn}\)（\(m \geqslant 1\)，\(n \geqslant 1\)）构成零邻域基本系的拓扑。）
\item
  把 Šmulian 定理推广到 Fréchet 空间的一个严格归纳极限空间（II, p.~33）。
\end{enumerate}

\begin{enumerate}
\def\labelenumi{\arabic{enumi})}
\setcounter{enumi}{2}
\tightlist
\item
  \begin{enumerate}
  \def\labelenumii{\alph{enumii})}
  \tightlist
  \item
    设 E 为分离的局部凸空间，\(E'\) 为其对偶；假设 \(E'\) 中存在对拓扑 \(\sigma(E', E)\) 处处稠密的可数集。证明 E 的拓扑（相应地，\(\sigma(E, E')\)）细于某个可度量化的局部凸拓扑。
  \end{enumerate}
\end{enumerate}

\begin{enumerate}
\def\labelenumi{\alph{enumi})}
\setcounter{enumi}{1}
\item
  设 \((x_{n})\) 为 E 中点的序列，使得从 \((x_{n})\) 中抽取的每个序列对初始拓扑（相应地，拓扑 \(\sigma (\mathrm{E},\mathrm{E}^{\prime})\)）都有极限点。证明存在从 \((x_{n})\) 中抽取的序列，它在 E 中对初始拓扑（相应地，对拓扑 \(\sigma (\mathrm{E},\mathrm{E}^{\prime})\)）收敛。（设 \((a_n^{\prime})\) 为 \(\mathbf{E}'\) 中对 \(\sigma (\mathrm{E}',\mathrm{E})\) 处处稠密的序列；抽取一个序列 \((y_{n})\) （取自 \((x_{n})\) ），使 \((\langle y_n,a_p'\rangle)\) 对每个指标 \(p\) 都趋于一个极限，并证明序列 \((y_{n})\) 对初始拓扑（相应地，对 \(\sigma (\mathrm{E},\mathrm{E}^{\prime}))\)）只有一个极限点。）
\item
  为使 E 的子集 A 对初始拓扑（相应地，对 \(\sigma(E, E')\)）相对紧，必要且充分的是：从 A 中点的每个序列 \((x_n)\) 都能抽出一个序列 \((x_{n_k})\) ，它对初始拓扑（相应地，对 \(\sigma(E, E')\)）收敛于 E 中的一点（利用练习 2, a)）。
\end{enumerate}

\begin{enumerate}
\def\labelenumi{\arabic{enumi})}
\setcounter{enumi}{3}
\item
  设 \(\mathrm{E}\) 为不满足第一可数公理（I, p.~25, 练习 1）的 Banach 空间 \(l^{\infty}(\mathbb{N})\)，\(\mathrm{E}'\) 为其对偶。对每个整数 \(n \geqslant 0\) ，设 \(e_n'\) 为 E 上的连续线性型，它把每个 \(x = (\xi_n) \in \mathrm{E}\) 对应到该序列的第 \(n\) 项。证明序列 \((e_n')\) 在 \(\mathrm{E}'\) 中对 \(\sigma(\mathrm{E}', \mathrm{E})\) 是完全的；此外，\((e_n')\) 的每个抽取序列在 \(\mathrm{E}'\) 中对拓扑 \(\sigma(\mathrm{E}', \mathrm{E})\) 都有极限点，但不存在 \((e_n')\) 的任何抽取序列在这个拓扑下于 \(\mathrm{E}'\) 中收敛。
\item
  \begin{enumerate}
  \def\labelenumii{\alph{enumii})}
  \tightlist
  \item
    设 X 为紧空间，H 为由 X 上连续数值函数组成的空间 \(\mathcal{C}(\mathrm{X})\) 的任意子集。设 \(f_{0}\) 为 \(\mathcal{C}(\mathrm{X})\) 中的一点，它对简单收敛拓扑 \(T_{s}\) 属于 H 在 \(\mathcal{C}(\mathrm{X})\) 中的闭包。证明存在 H 的可数子集 \(H_{0}\) ，使 \(f_{0}\) 属于 \(H_{0}\) 的闭包（对 \(T_{s}\) ）。（证明对每对整数 m \textgreater{} 0, n \textgreater{} 0，存在 H 的有限子集 \(\mathrm{H}(m, n)\) 具有下列性质：对 X 中 m 个点的每个组 \(t_{k}\)（\(1 \leqslant k \leqslant m\)），存在 \(f \in \mathrm{H}(m, n)\) 使 \(|f_{0}(t_{k}) - f(t_{k})| \leqslant 1/n\) 对 \(1 \leqslant k \leqslant m\) 成立。b) 设 E 为可度量化的局部凸空间，\(E'\) 为其对偶，H 为 E 的子集。证明若 \(x_{0}\) 对弱化拓扑 \(\sigma(\mathrm{E}, \mathrm{E}')\) 属于 H 的闭包，则存在 H 的可数子集 \(H_{0}\) ，使 \(x_{0}\) 对此拓扑属于 \(H_{0}\) 的闭包。（利用 a)，注意 \(E'\) 是对 \(\sigma(\mathrm{E}', \mathrm{E})\) 紧的集合所组成的可数族的并。）
  \end{enumerate}
\end{enumerate}

¶ 6) a) 设 X 为紧空间，H 为乘积空间 \(\mathbf{R}^{\mathrm{X}}\) 中由 X 上连续函数组成的凸子集。假设 H 中非空凸闭子集的每个递减序列都有非空的交。证明闭包 \(\overline{\mathrm{H}}\) （H 在 \(\mathbf{R}^{\mathrm{X}}\) 中的）是紧的且由 X 上的连续函数组成。（用反证法：考虑一个不连续的函数 \(u \in \overline{\mathrm{H}}\) ；证明此时将存在一点 \(a \in \mathrm{X}\) 、一个数 \(\delta > 0\) 、X 的点的一个序列 \((x_{n})\) 以及 H 中函数的一个序列 \((f_{m})\) ，使得：

\(1^{\circ}\left|u(x_{n}) - u(a)\right|\geqslant \delta\) 对所有 \(n;2^{\circ}\left|f_{m}(x_{n}) - f_{m}(a)\right|\leqslant \frac{\delta}{8}\) 对于 \(m\leqslant n;3^{\circ}\left|u(x_{n}) - f_{m}(x_{n})\right|\leqslant \frac{\delta}{8}\) 以及

\(\left|u(a) - f_m(a)\right| \leqslant \frac{\delta}{8}\) 对 \(m \geqslant n + 1\) 成立。设 \(b\) 为序列 \((x_n)\) 的一个极限点，并设 \(f\) 为属于诸 \(\mathbf{A}_m\) 之交的一个函数，其中 \(\mathbf{A}_m\) 是 \(\mathbf{H}\) 中由所有 \(f_k\)（\(k \geqslant m\)）组成的集合的闭凸包。

\begin{enumerate}
\def\labelenumi{\alph{enumi})}
\setcounter{enumi}{1}
\tightlist
\item
  设 E 为分离且拟完备的局部凸空间，\(E'\) 为其对偶。设 H 为 E 的凸子集，使得 H 的非空闭凸子集的每个递减序列都有非空的交；证明 H 在 E 中对拓扑 \(\sigma(\mathrm{E}, \mathrm{E}')\) 相对紧。（归结为 E 完备的情形；把 E 看作嵌入于 \(E'^{*}\) ，并利用 a) 以及 III, p.~21, 推论 1。）
\end{enumerate}

\begin{enumerate}
\def\labelenumi{\arabic{enumi})}
\setcounter{enumi}{6}
\item
  设 \(\mathbf{E}\) 为满足第一可数公理的 Fréchet 空间，\(\mathbf{E}'\) 为其对偶。证明：为使凸子集 \(\mathbf{A}'\) （\(\mathbf{E}'\) 的）对 \(\sigma(\mathbf{E}', \mathbf{E})\) 为闭，只需：若点的序列 \((x_n')\) （属于 \(\mathbf{A}'\) ）有一个极限 \(a'\) （在 \(\mathbf{E}'\) 中，对 \(\sigma(\mathbf{E}', \mathbf{E})\) ），则 \(a' \in \mathbf{A}'\) 。
\item
  设 \(\mathrm{F}\) 与 \(\mathrm{G}\) 为处于分离对偶中的两个向量空间。证明 IV, p.~52, 练习 6, a) 中的性质 \(\alpha)\) 与 \(\beta)\) 也等价于下列性质：
\end{enumerate}

\(\gamma)\) 赋有 \(\tau (\mathrm{F},\mathrm{G})\) 的 F 是拟完备的，且 F 中点的每个有界序列对 \(\sigma (\mathrm{F},\mathrm{G})\) 都有极限点（参见 IV, p.~35, 定理 1）；

δ) 赋有 \(\tau(\mathrm{F},\mathrm{G})\) 的 F 是拟完备的，且 F 中非空闭凸有界集的每个递减序列都有非空的交（参见练习 6, b)）。

\begin{enumerate}
\def\labelenumi{\arabic{enumi})}
\setcounter{enumi}{8}
\tightlist
\item
  设 E 为分离且拟完备的局部凸空间；为使 E 半自反，必要且充分的是：E 的每个含有可数处处稠密子集的闭向量子空间都是半自反的（参见 IV, p.~35, 定理 1）。
\end{enumerate}

¶ 10) 设 E 为分离、拟完备且非半自反的局部凸空间，H 为 E 中包含原点的一个闭超平面。设 \((\mathbf{C}_{n})\) 为非空、闭凸且有界之集的递减序列，这些集合包含于 H 中、不含 0 且其交为空（练习 8）。设 x 为不属于 H 的一点，并设 A 为

诸集合 \(\left(1 - \frac{1}{n}\right)x + C_n\)（\(n > 0\)）的并集的闭凸均衡包。

\begin{enumerate}
\def\labelenumi{\alph{enumi})}
\item
  证明不存在 A 的平行于 H 的支撑超平面（对 \(y \in x + H\) ，注意存在整数 n 使 \(y \notin x + C_{n}\) ）。
\item
  设 \(z \in \mathbf{H}\) 满足 \(z \notin \mathbf{C}_1\) ；证明两个闭凸有界集 \(\mathbf{A}\) 与 \(\mathbf{B} = x + z + \mathbf{C}_1\) 的并集的凸包不是闭的（证明 \(x + z\) 属于该凸包的闭包，但不属于该凸包）。
\end{enumerate}

\begin{enumerate}
\def\labelenumi{\arabic{enumi})}
\setcounter{enumi}{10}
\item
  设 A 为无限集，\(E = \overline{\mathcal{H}}(A)\)（IV, p.~47, 练习 1），\(E' = \ell^{1}(A)\) 为其对偶，\(E'' = \ell^{\infty}(A)\) 为其双对偶（IV, p.~54, 练习 14）。证明 \(E'\) 的每个对 \(\sigma(E', E'')\) 紧的子集都是强紧的（利用 Šmulian 定理以及 IV, p.~54, 练习 15）。
\item
  设 E 为非自反的 Banach 空间。证明 E 中存在一个闭的、非自反的且余维数无限的向量子空间 M。（设 \((x_{n})\) 为 E 中对拓扑 \(\sigma(\mathrm{E},\mathrm{E}')\) 没有极限点的有界序列（IV, p.~68, 练习 8）；用归纳法构造一个序列 \((x_{n_{k}})\) （自 \((x_{n})\) 中抽取的）和一个拓扑自由序列 \((y_{k})\) ，使 \(\| x_{n_{k}} - y_{k} \| \leqslant 1/k\) 对 \(k \geqslant 1\) 成立，并考虑由诸 \(y_{2k}\) 生成的 E 的闭向量子空间。
\end{enumerate}

¶ 13) 设 E 为满足第一可数公理的非自反 Banach 空间，M 为 E 中余维数无限的非自反闭向量子空间（练习 12）。设 \((x_{n})\) 为单位球面中处处稠密的序列，C 为序列 \((x_{n}/n)\) 的闭凸均衡包；C 在 E 中强紧，且 \(C + M = A\) 是闭凸集（GT, III, § 4, No.~1, 命题 1 的推论 1）。设 S 为 E 中的单位球，\(B = A \cap S\) 。

\begin{enumerate}
\def\labelenumi{\alph{enumi})}
\item
  证明 0 不是 A 的内点，并由此推出存在 \(x_0 \in E\) 使 \(\lambda x_0 \notin A\) 对 \(\lambda > 0\) 成立。
\item
  证明不存在过 0 的 B 的支撑超平面（注意这样的超平面必是 C 的支撑超平面）。
\item
  设 \(\mathrm{U}_0 = \mathrm{M} \cap \mathrm{S}\) ，\((\mathrm{U}_n)_{n \geqslant 1}\) 为闭凸、有界且非空之集的递减序列，满足 \(\mathrm{U}_1 \subset \frac{1}{2}\mathrm{U}_0\) 且诸 \(\mathrm{U}_n\) 的交为空（IV, p.~68, 练习 8）。设 \(\mathbf{F}\) 为诸集合 \(\frac{1}{n} x_0 + \mathrm{U}_n\)（\(n \geqslant 1\)）的并集的闭凸包。证明 \(\mathbf{B} \cap \mathbf{F} = \emptyset\) ，但不存在分离 \(\mathbf{B}\) 与 \(\mathbf{F}\) 的闭超平面（利用 \(b\) )）。
\end{enumerate}

\begin{enumerate}
\def\labelenumi{\arabic{enumi})}
\setcounter{enumi}{13}
\tightlist
\item
  设 \(\mathbf{E}\) 为满足第一可数公理的 Banach 空间。元素的序列 \((e_n)_{n\geqslant 0}\) （\(\mathbf{E}\) 中的）称为 Banach 基，如果下列性质成立：对每个 \(x\in \mathbb{E}\) ，存在唯一的标量序列 \((a_{n})\) 使 \(x = \sum_{n = 0}^{\infty}a_{n}e_{n}\) ，其中右端的级数
\end{enumerate}

收敛。

\begin{enumerate}
\def\labelenumi{\alph{enumi})}
\tightlist
\item
  证明族 \((e_n)\) 是完全的且自由的。设 \(E_n\) 为由诸 \(e_m\) （指标满足 \(m \leqslant n\) ）生成的 E 的（闭）向量子空间，\(P_n\) 为 E 到 \(E_n\) 上的由 \(P_n \cdot (\sum_{m=0}^{\infty} a_m e_m) = \sum_{m=0}^{n} a_m e_m\) 定义的投影算子。证明诸 \(P_n\) 是连续线性映射，且
\end{enumerate}

\(\sup_{n}\| P_{n}\| < + \infty .\)（考虑 E 上由 \(\left\| \sum_{n = 0}^{\infty}a_n e_n\right\| = \sup_{n}\left\| \sum_{m = 0}^{n}a_me_m\right\|\)

证明 E 对此范数完备，并由此推出它与 E 上给定的范数等价（参见 I, p.~17, 定理 1）。）证明对每对整数 \(p < q\) ，投影 \(P_{q,p}\)（从 \(\mathrm{E}_q\) 到 \(\mathrm{E}_p\) 、平行于由诸 \(e_m\) （指标 \(m\) 满足 \(p + 1 \leqslant m \leqslant q\) ）生成的向量子空间）的范数被一个与 \(p, q\) 无关的数所界定。

\begin{enumerate}
\def\labelenumi{\alph{enumi})}
\setcounter{enumi}{1}
\tightlist
\item
  反之，设 \((e_{n})\) 为 E 中元素的一个完全序列，它是一个自由族，且诸投影 \(P_{q,p}\) （对 \(0 \leqslant p < q\) ）的范数被一个与 p 和 q 无关的数 M 所界定。证明 \((e_{n})\) 是 E 的 Banach 基。（先证明序列 \((e_{n})\) 拓扑自由并定义诸投影算子 \(P_{n}\) ；于是对所有 n 有 \(\|P_{n}\| \leqslant M\) 。其次注意，若 \(d(x, E_{n})\) 是点 \(x \in E\) 到 \(E_{n}\) 的距离，则 \(\|x - P_{n} \cdot x\| \leqslant (M + 1) d(x, E_{n})\) 。) \(^{1}\)
\end{enumerate}

\(^{1}\) 显然，E 中存在 Banach 基蕴含 E 满足第一可数公理。但也存在满足第一可数公理而没有 Banach 基的 Banach 空间的例子（P. ENFLO, Acta Math., t. CXXX (1973), p.~309-317）。

¶ 15) 设 E 为具有 Banach 基 \((e_n)\) 的 Banach 空间（练习 14）；则存在唯一的序列 \((e_n')\) 在 E 的对偶 \(E'\) 中，使每个 \(x \in E\) 关于 Banach 基 \((e_n)\) 的表达式可写成 \(x = \sum_{n=0}^{\infty} \langle x, e_n' \rangle e_n\) 。

\begin{enumerate}
\def\labelenumi{\alph{enumi})}
\item
  设 \(\mathrm{F}_n\) 为 \(\mathbf{E}\) 的由诸 \(e_m\) （指标满足 \(m \geqslant n\) ）生成的闭向量子空间。对每个 \(x' \in \mathbf{E}'\) ，用 \(\| x' \|_n\) 表示线性型 \(x'\) 在 \(\mathrm{F}_n\) 上的限制的范数。证明：为使 \((e_n')\) 是 \(\mathbf{E}'\) 的 Banach 基，必要且充分的是对每个 \(x' \in \mathbf{E}'\) ，序列 \((\| x' \|_n)\) 趋于 0。（考虑转置 \({}^t P_n\) 并估计范数 \(\| {}^t P_n \cdot x' - x' \|\) 。）这时称 Banach 基 \((e_n)\) 为收缩的。
\item
  假设 Banach 基 \((e_n)\) 是收缩的。证明对每个点 \(x''\) （E 的双对偶 \(\mathrm{E}''\) 的），和的序列 \(\sum_{m=0}^{n} \langle e_m', x'' \rangle e_m\) 在 E 中有界（考虑转置 \(t^t P_n\) ）
\end{enumerate}

在 \(\mathrm{E}''\) 中）。反之，对每个标量序列 \((a_{n})\) ，若和的序列 \(\sum_{m=0}^{n} a_{m} e_{m}\)

在 E 中有界，则存在唯一的 \(x'' \in E''\) 使 \(\langle e_{n}', x'' \rangle = a_{n}\) 对所有 n 成立（利用 \(E''\) 中闭球关于拓扑 \(\sigma(E'', E')\) 的紧性）。

\begin{enumerate}
\def\labelenumi{\alph{enumi})}
\setcounter{enumi}{2}
\item
  E 中的 Banach 基 \((e_n)\) 称为完全的，如果对每个标量序列 \((a_n)\) ，当和的序列 \(\sum_{m=0}^{n} a_m e_m\) 有界时，级数 \(\sum_{n=0}^{\infty} a_n e_n\) 收敛。若 \((e_n)\) 是 E 的收缩基，则基 \((e_{n}^{\prime})\) （\(E^{\prime}\) 的）是完全的（利用 \(E^{\prime}\) 中闭球关于 \(\sigma(\mathrm{E}^{\prime},\mathrm{E})\) 的紧性）。 d) 一般地，序列 \((e_{n}^{\prime})\) 是闭子空间 \(F^{\prime}\) （即 E 的强对偶 \(E^{\prime}\) 中由诸 \(e_{n}^{\prime}\) 生成的闭子空间）的 Banach 基，且存在一个从 E 到强对偶 \(F^{\prime\prime}\) （\(F^{\prime}\) 的）中的单射连续线性映射 J，使 \(\langle J.x,z^{\prime}\rangle=\langle x,z^{\prime}\rangle\) 对所有 \(x\in E\) 和所有 \(z^{\prime}\in F^{\prime}\) 成立。证明存在常数 K\textgreater0 使 \(\|J.x\|\geqslant K.\|x\|\) ，且若基 \((e_{n})\) 是完全的，则 J 是从 E 到拓扑向量空间 \(F^{\prime\prime}\) 上的同构。
\item
  证明为使 \(\mathbf{E}\) 自反，必要且充分的是基 \((e_n)\) 既收缩又完全（利用 \(b\) ）。
\end{enumerate}

\begin{enumerate}
\def\labelenumi{\arabic{enumi})}
\setcounter{enumi}{15}
\tightlist
\item
  \begin{enumerate}
  \def\labelenumii{\alph{enumii})}
  \tightlist
  \item
    设 \(\mathbf{E}\) 为 Banach 空间。对点的无穷序列 \((x_{n})\) （\(\mathbf{E}\) 中的），下列性质等价：
  \end{enumerate}
\end{enumerate}

\(\alpha)\) 通项为 \(x_{n}\) 的级数交换收敛（GT, III, § 5, No.~7）。

\(\beta)\) 对 \(\mathbf{N}\) 的每个子集 I ，由序列 \((x_{n})_{n\in \mathrm{I}}\) 定义的级数收敛（GT, III, § 5, No.~3, 命题 2 与 § 5, 练习 4）。

\(\gamma)\) 对每个数的序列 \((\varepsilon_{n})\) （每项等于 1 或 \(-1\) ），通项为 \(\varepsilon_{n}x_{n}\) 的级数收敛。

\(\delta)\) 对每个 \(\varepsilon > 0\) ，存在有限子集 \(J\) （\(N\) 的），使对每个有限子集 \(H\) （\(N\) 的，与 \(J\) 不相交的），有 \(\| \sum x_n \| \leqslant \varepsilon\) 。

\begin{enumerate}
\def\labelenumi{\alph{enumi})}
\setcounter{enumi}{1}
\tightlist
\item
  设 \((e_n)\) 为 E 的 Banach 基。下列性质等价：
\end{enumerate}

\(\alpha)\) 对每个置换 \(\pi\) （\(\mathbf{N}\) 的），序列 \((e_{\pi(n)})\) 是 E 的 Banach 基。

\(\beta)\) 对每项等于 1 或 -1 的数的序列 \((\varepsilon_{n})\) ，序列 \((\varepsilon_{n}e_{n})\) 是 E 的 Banach 基。

\(\gamma)\) 对 E 中每个 \(x = \sum_{n=0}^{\infty} \xi_n e_n\) 以及每个序列 \((\eta_n)_{n \in \mathbb{N}}\) （满足 \(|\eta_n| \leqslant |\xi_n|\) 对所有 \(n\) 成立），通项为 \(\eta_n e_n\) 的级数在 E 中收敛。

δ) 对 E 中每个 \(x = \sum_{n=0}^{\infty} \xi_n e_n\) 以及每个严格递增序列 \((n_k)_{k \in \mathbf{N}}\) （由整数 \(\geqslant 0\) 组成的），通项为 \(\xi_{n_k} e_{n_k}\) 的级数在 E 中收敛。

ε) 存在实数 M \textgreater{} 0 ，使对 N 的每个有限子集 J 以及 E 中每个 \(x = \sum_{n=0}^{\infty} \xi_{n} e_{n}\) ，有 \(\left\|\sum_{n \in J} \xi_{n} e_{n}\right\| \leqslant M \|x\|\) 。

（为证明 \(\alpha\) ) 蕴含 \(\varepsilon\) )，按 IV, p.~69, 练习 14 的方式论证。为证明 \(\beta\) 蕴含 \(\gamma\) )，归结到 \(\xi_{n}\) 与 \(\eta_{n}\) 为实数的情形，并考虑和 \(\sum_{n=p}^{q} \langle \eta_{n} e_{n}, x' \rangle\) （对 \(x' \in E'\) ）。）

当这些条件满足时，称 \((e_n)\) 为 E 的无条件 Banach 基。c) 假设 \((e_n)\) 是无条件基。证明存在实数 \(K > 0\) ，使对每项等于 1 或 -1 的每个序列 \((\varepsilon_n)\) 以及 E 中所有 \(x = \sum_{n=0}^{\infty} \xi_n e_n\) ，有 \(\| \sum_{n=0}^{\infty} \varepsilon_n \xi_n e_n \| \leqslant M \| x\|\) （方法同 IV, p.~69, 练习 14）。由此推出，对每个

有界的标量序列 \((\lambda_{n})\) 以及 X 中所有 \(x = \sum_{n=0}^{\infty} \xi_{n} e_{n}\) ，有

\[
\left\| \sum_ {n = 0} ^ {\infty} \lambda_ {n} \xi_ {n} e _ {n} \right\| \leqslant 2 K \sup _ {n} | \lambda_ {n} |. \left\| \sum_ {n = 0} ^ {\infty} \xi_ {n} e _ {n} \right\|
\]

（论证方式同 \(b\) ) 中证明 \(\beta\) 蕴含 \(\gamma\) 的做法）。

¶ 17) a) 设 E 为具有无条件 Banach 基 ( \(e_n\) ) 的 Banach 空间（练习 16）。证明若基 ( \(e_n\) ) 不是收缩的（IV, p.~70, 练习 15），则存在数 \(\alpha > 0\) 、线性型 \(x' \in E'\) （满足 \(\|x'\| = 1\) ）、整数的严格递增序列 ( \(n_k\) ) ，并且对每个 \(k\) ，存在元素 \(y_k\) ，它是诸 \(e_j\) （指标满足 \(n_k \leqslant j \leqslant n_{k+1}\) ）的线性组合，且满足 \(\|y_k\| \leqslant 1\) 与 \(\langle y_k, x' \rangle \geqslant \alpha\) 。由此推出，对每个标量的有限序列 ( \(\lambda_j\) ) \(_{1 \leqslant j \leqslant m}\) ，有 \(\|\sum_{j=1}^{m} \lambda_j y_j\| \geqslant \frac{\alpha}{2K} \sum_{j=1}^{m} |\lambda_j|\) （利用练习 16, c)）。最后得出存在从 \(\ell^1 (\mathbf{N})\) 到 E 的一个闭子空间上的拓扑向量空间同构。

\begin{enumerate}
\def\labelenumi{\alph{enumi})}
\setcounter{enumi}{1}
\item
  由 \(a\) ) 推出：若强对偶 \(\mathrm{E}'\) （\(\mathrm{E}\) 的）满足第一可数公理，则每个无条件基 \((e_n)\) （\(\mathrm{E}\) 的）都是收缩的，从而 \((e_n')\) 是 \(\mathrm{E}'\) 的无条件基（IV, p.~70, 练习 15, a)）。（注意，若 \(\mathrm{E}\) 的某个闭向量子空间同构于 \(\ell^1(\mathbf{N})\) ，则 \(\mathrm{E}'\) 不可能满足第一可数公理。）
\item
  证明若 E 具有无条件基，且 E 的强双对偶 \(E''\) 满足第一可数公理，则 E 自反。（利用 IV, p.~51, 练习 25，注意 E 的强对偶 \(E'\) 满足第一可数公理；然后用 IV, p.~70, 练习 15, c) 与 IV, p.~53, 练习 11。）
\end{enumerate}

¶ 18) a) 在由实数的全体无穷序列组成的空间 \(\mathbf{R}^{\mathbf{N}}\) 中，考虑由一切序列 \(x = (\xi_n)\) （对这些序列，数

\[
\| x \| = \sup \left(\left(\xi_ {p _ {1}} - \xi_ {p _ {2}}\right) ^ {2} + \left(\xi_ {p _ {2}} - \xi_ {p _ {3}}\right) ^ {2} + \dots + \left(\xi_ {p _ {m - 1}} - \xi_ {p _ {m}}\right) ^ {2} + \xi_ {p _ {m + 1}} ^ {2}\right) ^ {1 / 2}
\]

为有限）所组成的集 J，其中上确界取遍所有整数 \(m \geqslant 1\) 与整数的所有严格递增序列 \(p_{1} < p_{2} < \cdots < p_{m+1}\) 。证明 \(\|x\|\) 是 J 上的一个范数，且 J 对此范数是 Banach 空间。对每个 \(x \in E\) ，证明序列 \((\xi_{n})\) 有有限极限 \(u(x)\) ，且 u 是 E 上非零的连续线性型；用 \(J_{0}\) 表示 J 中由方程 \(u(x) = 0\) 给定的闭超平面（R. C. James 空间）。

\begin{enumerate}
\def\labelenumi{\alph{enumi})}
\setcounter{enumi}{1}
\item
  证明向量 \(e_n = (\delta_{mn})_{m \geq 0}\) 构成 \(J_0\) 的 Banach 基，且此基是收缩的（IV, p.~70, 练习 15）。（为证明后一点，按练习 17, a) 的方式论证，证明所构造的序列 \((y_k)\) 使通项为 \(y_k / k\) 的级数在 E 中收敛；这将导出矛盾。）
\item
  证明 \(\mathrm{J}_0\) 到自身的恒等映射可以延拓为从强双对偶 \(\mathrm{J}_0''\) 到 \(\mathrm{J}\) 上的拓扑向量空间同构，且使 \(\mathrm{J}_0'' / \mathrm{J}_0\) 是 1 维的（利用 IV, p.~70, 练习 15, b)）。由此推出 \(\mathrm{J}_0\) 的任何 Banach 基都不可能是无条件的（练习 17, c)）。
\end{enumerate}

* \(d\) ) 设 \(\mathrm{H}_1,\mathrm{H}_2\) 为 \(\mathbf{J}_0\) 中分别由诸 \(e_{2n}\) 与诸 \(e_{2n + 1}\) （\(n\geqslant 0\) ）生成的两个闭向量子空间。证明作为拓扑向量空间，\(\mathrm{H}_{1}\) 与 \(\mathrm{H}_{2}\) 同构于 Hilbert 空间 \(\ell^2 (\mathbf{N})\) ，并且 \(\mathbf{J}_0\) 不是 \(\mathrm{H}_1\) 与 \(\mathrm{H}_2.\) 的和。

\begin{enumerate}
\def\labelenumi{\alph{enumi})}
\setcounter{enumi}{4}
\tightlist
\item
  证明在 \(\mathrm{J}_0\) 上不存在以 \(\mathrm{J}_0\) 的实局部凸空间结构为底层结构的复局部凸空间结构（参见~IV, p.~52, 练习 3）。
\end{enumerate}

